\documentclass[11pt]{article}
\usepackage[margin=0.82in]{geometry}
\usepackage{amsmath}
\usepackage{booktabs}
\usepackage{array}
\usepackage{float}
\usepackage{hyperref}
\hypersetup{colorlinks=true,urlcolor=blue,citecolor=blue,linkcolor=blue}

\title{Project Four: Fund X Factor Model Analysis}
\author{Prepared by Harman Singh for Professor Phelim Boyle}
\date{July 2026}

\begin{document}
\maketitle

\section*{Data and Method}
The analysis uses Fund X's net-of-fees monthly returns from January 2017 through
December 2024 and the official monthly q5 factor file from
\url{https://global-q.org/factors.html}. The 96 observations are merged by
year and month. Every factor and return is measured in percent per month. The
dependent variable is Fund X's excess return, $r_{X,t}-r_{f,t}$, where
$r_{f,t}$ is the \texttt{R\_F} observation supplied in the q5 file.

Three time-series regressions are estimated: the CAPM with $R_{MKT}$; the
Hou--Xue--Zhang q-factor model with $R_{MKT}$, $R_{ME}$, $R_{IA}$, and
$R_{ROE}$; and the augmented q5 model, which adds $R_{EG}$. Coefficients are
ordinary least squares estimates. The reported $t$-statistics use Newey--West
heteroskedasticity- and autocorrelation-consistent covariance estimates with
the assignment-mandated Bartlett lag length $L=4$.

\clearpage
\section*{Table 1: Full Regression Results}
\begin{table}[H]
\centering
\caption{Fund X Factor Regressions: CAPM, q-Factor, and q5}
\label{tab:factor-regressions}
\small
\begin{tabular}{lccc}
\toprule
 & Model 1 & Model 2 & Model 3 \\
 & CAPM & q-factor & q5 \\
\midrule
$\alpha$ (monthly, \%) &
3.015*** & 3.041*** & 2.965*** \\
 & [4.93] & [5.24] & [5.48] \\[3pt]
$\beta_{MKT}$ &
-0.056 & -0.022 & 0.010 \\
 & [-0.50] & [-0.26] & [0.14] \\[3pt]
$\beta_{ME}$ (size) &
-- & -0.222 & -0.093 \\
 & -- & [-1.29] & [-0.69] \\[3pt]
$\beta_{IA}$ (investment) &
-- & 0.122 & 0.273 \\
 & -- & [1.12] & [1.39] \\[3pt]
$\beta_{ROE}$ (profitability) &
-- & -0.161 & -0.350 \\
 & -- & [-1.26] & [-1.88] \\[3pt]
$\beta_{EG}$ (expected growth) &
-- & -- & 0.366 \\
 & -- & -- & [1.07] \\
\midrule
$R^2$ & 0.006 & 0.040 & 0.069 \\
Adjusted $R^2$ & -0.005 & -0.002 & 0.017 \\
$N$ & 96 & 96 & 96 \\
\bottomrule
\end{tabular}

\vspace{0.2cm}
\begin{minipage}{0.96\linewidth}
\footnotesize
\textit{Note:} Sample: January 2017 -- December 2024 (N = 96 monthly observations).
Dependent variable: Fund X monthly excess return ($r_X-r_f$). q5 factor returns
($R_{MKT}$, $R_{ME}$, $R_{IA}$, $R_{ROE}$, $R_{EG}$) are from Hou, Mo,
Xue and Zhang (2021) via global-q.org. Standard errors are Newey--West HAC with
$L=4$ lags (Newey and West, 1987); $t$-statistics in brackets. No coefficient is
asterisked unless its $|t|\geq1.96$. *** denotes $p<0.001$.
\end{minipage}
\end{table}

\clearpage
\section*{Table 2: Alpha Summary}
\begin{table}[H]
\centering
\caption{Summary of Alpha Estimates Across Factor Models}
\label{tab:alpha-summary}
\small
\begin{tabular}{lrrrr}
\toprule
Model & Monthly $\alpha$ & NW $t$-stat & Annual $\alpha$ & $R^2$ \\
\midrule
CAPM (1-factor) & 3.015*** & 4.93 & 42.82\% & 0.006 \\
q-factor (4-factor) & 3.041*** & 5.24 & 43.27\% & 0.040 \\
q5 (5-factor) & 2.965*** & 5.48 & 42.00\% & 0.069 \\
\bottomrule
\end{tabular}

\vspace{0.2cm}
\begin{minipage}{0.94\linewidth}
\footnotesize
\textit{Note:} *** $p<0.001$. Annual alpha computed as
$(1+\text{monthly\_alpha}/100)^{12}-1$. Sample: January 2017 -- December 2024,
$N=96$. NW $t$-statistics use $L=4$ lags.
\end{minipage}
\end{table}

\section*{Results}
Across all three specifications, Fund X retains a large positive intercept. The CAPM
monthly alpha is 3.01\% with a Newey--West $t$-statistic of
4.93. The q-factor estimate is 3.04\%
($t=5.24$), and the q5 estimate is 2.97\%
($t=5.48$). The full range across specifications is only
7.6 basis points, while the change from CAPM to q5 is
-5.0 basis points. The corresponding compounded annual alpha estimates
range from 42.00\% to 43.27\%. All three
intercepts have $p<0.001$. The $t$-statistic rises as factors are added, despite the
slight decline in the q5 point estimate, because the HAC standard error becomes smaller.
The additional factors therefore do not absorb Fund X's average excess return.

The market loading is economically small and statistically insignificant in every
model: it moves from -0.056 in the CAPM to
0.010 in q5. In the q5 model, the size loading is
-0.093 ($t=-0.69$), the investment
loading is 0.273 ($t=1.39$), and the
expected-growth loading is 0.366
($t=1.07$). None is significant at the 5\% level. The
profitability loading, -0.350 with
$t=-1.88$, comes closest to conventional significance. If a
negative loading were stable and statistically reliable, it would mean that Fund X
tends to move against the high-minus-low ROE factor rather than earning its returns by
harvesting the profitability premium. Here the evidence is suggestive, not conclusive.

The CAPM explains only 0.6\% of monthly variation, compared with
4.0\% for the q-factor model and 6.9\% for
q5. Adjusted $R^2$ values are -0.005,
-0.002, and 0.017, respectively. The
negative adjusted values in the first two models are valid: after penalising additional
parameters, those specifications do not improve on an intercept-only benchmark. These
figures are far below the 85--95\% market-model $R^2$ often seen for diversified equity
funds. The q factors therefore explain little of Fund X's month-to-month return
variation. Low explanatory power is not itself proof of skill, but together with the
large intercept it shows that these standard linear equity factors do not account for
the reported performance.

The main conclusion is that Fund X's alpha survives both the q-factor and q5 controls.
Within this sample and model class, its excess return remains largely unexplained by
known systematic equity-factor premia.

\clearpage
\section*{Assumption-Aware Diagnostic Assessment}
The required $L=4$ estimates are the primary results. The literal bandwidth expression
$4(T/100)^{2/9}$ equals 3.9639 at $T=96$ and
would floor to 3, so $L=3$ is also checked.
The q5 alpha $t$-statistic is 5.78 at $L=3$,
5.48 at $L=4$, 5.07 at $L=6$, and
4.34 at $L=12$. The conclusion is not driven by the lag choice.

The q5 residuals remain serially dependent: the Ljung--Box statistic through 12 lags is
51.38 with $p<0.001$. They are also strongly non-normal
(Jarque--Bera $p<0.001$), and the Breusch--Pagan test is borderline
($p=0.058$). These findings justify HAC
inference and caution against relying on textbook iid OLS standard errors. A
6-month circular block-pairs bootstrap, which preserves local
dependence and the joint factor-return observations, gives a 95\% q5 alpha interval of
[1.85\%, 4.10\%].

March 2020 is the most influential observation (Cook's distance
1.28). Excluding it leaves a q5 alpha of
2.77\% with $t=5.57$.
Across all 96 leave-one-out regressions, alpha ranges only from
2.77\% to
3.04\%.

Parameter stability is the main qualification. A CUSUM test rejects a constant q5
coefficient vector ($p=0.0002$). In an objective midpoint
split, q5 alpha is 3.92\%
($t=6.01$) in 2017--2020 and
1.74\%
($t=3.31$) in 2021--2024. The alpha remains positive,
but its level is not stable. Newey--West inference cannot correct changing coefficients,
omitted factors, nonlinear exposures, or a mismatch between a US equity factor model and
the fund's underlying strategy. The result should therefore be stated as a large return
unexplained by these models, not as proof that every possible risk exposure has been ruled out.

\section*{References}
\begin{description}
\item Hou, K., Xue, C., and Zhang, L. (2015).
Digesting anomalies: An investment approach.
\textit{Review of Financial Studies}, 28(3), 650--705.
\item Hou, K., Mo, H., Xue, C., and Zhang, L. (2021).
An augmented q-factor model with expected growth.
\textit{Review of Finance}, 25(1), 1--41.
\item Newey, W. K. and West, K. D. (1987).
A simple, positive semi-definite, heteroskedasticity and autocorrelation consistent covariance matrix.
\textit{Econometrica}, 55(3), 703--708.
\end{description}

\end{document}
